\section{Theory}
\label{sec:theory}

Let \(\mathcal A\) and \(\mathcal U\) denote Banach spaces of input and output fields. A PDE family induces a solution operator \(\mathcal S:\mathcal A\to\mathcal U\). Let \(G\) be a Lie group acting on the full problem family through transformations \(T_g^{\mathcal A}\) and \(T_g^{\mathcal U}\). Closure of the full problem class under \(G\) means that the equation, coefficients, forcing terms, boundary or initial conditions, domain class, and observation map transform consistently. Under this assumption the exact solution operator is equivariant: \(\mathcal S(T_g^{\mathcal A}a)=T_g^{\mathcal U}\mathcal S(a)\). For a learned operator \(\mathcal G_\theta\), define the operator-level equivariance defect \(\Delta_\theta(a,g)=\mathcal G_\theta(T_g^{\mathcal A}a)-T_g^{\mathcal U}\mathcal G_\theta(a)\). Our regularizer samples \(g=\exp(\epsilon X_i)\) and penalizes
\begin{equation}
\mathcal L_{\mathrm{orb}}(\theta)=\mathbb E_{a,i,\epsilon}\left[\frac{\|\Delta_\theta(a,\exp(\epsilon X_i))\|_{\mathcal U}^2}{\epsilon^2+\eta}\right].
\end{equation}
The objective is label-free because both terms are model predictions transformed by known group actions.

\paragraph{Infinitesimal interpretation.} Assume \(T_{\exp(\epsilon X)}^{\mathcal A}\), \(T_{\exp(\epsilon X)}^{\mathcal U}\), and \(\mathcal G_\theta\) are differentiable at \(\epsilon=0\). Let \(X^{\mathcal A}a\) and \(X^{\mathcal U}u\) denote the induced infinitesimal actions. Taylor expansion gives
\begin{equation}
\mathcal G_\theta(T_{\exp(\epsilon X)}^{\mathcal A}a)-T_{\exp(\epsilon X)}^{\mathcal U}\mathcal G_\theta(a)=\epsilon\left(D\mathcal G_\theta(a)[X^{\mathcal A}a]-X^{\mathcal U}\mathcal G_\theta(a)\right)+O(\epsilon^2).
\end{equation}
Thus \(\|\Delta_\theta(a,\exp(\epsilon X))\|_{\mathcal U}^2/\epsilon^2\) converges to \(\|D\mathcal G_\theta(a)[X^{\mathcal A}a]-X^{\mathcal U}\mathcal G_\theta(a)\|_{\mathcal U}^2\) as \(\epsilon\to0\). The finite-orbit objective is therefore a stochastic finite-difference estimator of the squared Lie-algebraic equivariance defect, while avoiding explicit Jacobian-vector products.

\paragraph{Orbit error decomposition.} If the true solution operator is equivariant and \(T_g^{\mathcal U}\) is bounded, then for any \(a\in\mathcal A\) and \(g\in G\),
\begin{align}
\|\mathcal G_\theta(T_g^{\mathcal A}a)-\mathcal S(T_g^{\mathcal A}a)\|_{\mathcal U}
&\le \|\mathcal G_\theta(T_g^{\mathcal A}a)-T_g^{\mathcal U}\mathcal G_\theta(a)\|_{\mathcal U} \nonumber\\
&\quad+\|T_g^{\mathcal U}\|\,\|\mathcal G_\theta(a)-\mathcal S(a)\|_{\mathcal U}.
\end{align}
The first term is the learned equivariance defect and the second term is the in-distribution prediction error transported by the output action. This inequality motivates reducing \(\mathcal L_{\mathrm{orb}}\) when test inputs differ from training inputs by known PDE symmetries.

\paragraph{Approximate symmetries and cost.} For partial or approximate symmetries, replace the output norm by a weighted norm \(\|v\|_w^2=\|w_{a,g}\odot v\|^2\), where \(w_{a,g}\) can encode valid overlap after transformations, boundary exclusions, observation masks, or regions where transformed coefficients and forcing terms remain physically meaningful. With one sampled transform per minibatch, training uses one ordinary forward pass, one transformed forward pass, and a known output transformation. Inference is exactly the base map \(\mathcal G_\theta(a)\), so deployment latency and memory are unchanged relative to the underlying neural operator.
